% Grading criteria for CHEM 5260 Midterm Project
% Compile with: pdflatex criteria.tex
\documentclass[12pt,letterpaper]{article}

\usepackage[margin=1in]{geometry}
\usepackage{amsmath,amssymb}
\usepackage{braket}
\usepackage{parskip}

\begin{document}

\begin{center}
  {\large\bfseries CHEM 5260 Midterm Project --- Criteria for a Successful Report}\\[0.5em]
  {\normalsize Spring 2026}
\end{center}

\bigskip

A successful project report should function as a short scientific paper: it
should state a clear physical question, define the model precisely, derive
the relevant analytical benchmarks, present numerical results that test those
benchmarks, and explain the physical meaning of the observed agreement and
breakdown.  The following criteria define what such a report should contain.

\begin{enumerate}

\item \textbf{Organized structure.}
  The report should follow the format of a scientific paper and include all
  standard sections:
  \begin{itemize}
    \item \textit{Introduction} --- motivate the problem, place it in the
          context of quantum relaxation, Fermi's golden rule, and the
          wide-band approximation, and state clearly what the project
          investigates.
    \item \textit{Model and Methods} --- define the Hamiltonian(s), specify
          initial conditions and observables, and explain how the time
          evolution is computed.  The report should explain both the exact
          numerical propagation (e.g.\ exact diagonalization of the TDSE) and
          the analytical theory used as a benchmark.
    \item \textit{Results and Discussion} --- present numerical findings with
          figures, compare them to analytical expectations, and interpret
          them physically.
    \item \textit{Conclusion} --- summarize the key physical takeaways,
          limitations of the present models, and reasonable directions for
          future extensions.
  \end{itemize}

\item \textbf{Analytical content.}
  A successful report must contain a coherent analytical framework, not only
  numerical plots.  At minimum, the theory should establish the following:
  \begin{itemize}
    \item For \textit{Model A}, the Green's-function / self-energy derivation
          of exponential decay under the wide-band approximation, including
          the condition
          $\rho^{-1}\lesssim \Gamma \ll \Delta E$ and the FGR rate
          $\Gamma = 2\pi |V|^2 \rho$.
    \item For \textit{Model B}, the reduction of the $\{\ket{0},\ket{1}\}$
          subspace to a damped two-level problem when the doorway state
          $\ket{1}$ is assumed to undergo Markovian decay under the same
          wide-band approximation.
    \item The identification of overdamped, critically damped, and
          underdamped regimes in Model~B, with the threshold
          $V_{01}=\Gamma/4$.
    \item A clear statement of what assumptions enter the theory and where
          they may fail.
  \end{itemize}

\item \textbf{Numerical coverage.}
  The numerical investigation should test the analytical picture over the
  main control parameters of the model.  At minimum, the report should
  include and discuss:
  \begin{itemize}
    \item \textit{Model A}: dependence on bath size $N$, bath bandwidth
          $\Delta E$, and coupling strength $V_{1k}$;
    \item the connection between $N$, the recurrence time
          $T_\mathrm{rec}\sim 2\pi N/\Delta E$, and the finite-time validity
          of exponential decay;
    \item the role of the coupling disorder, including random-sign and
          Gaussian-random $V_{1k}$;
    \item \textit{Model B}: dependence on the upstream coupling $V_{01}$ and
          explicit demonstration of overdamped, critical, and underdamped
          transfer;
    \item the effect of detuning $E_0-E_1$ and its relation to the resonance
          / Lorentzian line shape;
    \item at least one explicit comparison between exact numerical results
          and the analytical benchmark wherever the benchmark is expected to
          apply.
  \end{itemize}

\item \textbf{Figure quality and interpretation.}
  The report should present figures at publication quality.  Each figure
  should have:
  \begin{itemize}
    \item a descriptive caption that explains what is shown and what
          parameters are used;
    \item properly labeled axes with the quantity name and unit (or an
          explicit statement that the quantity is dimensionless);
    \item legends or line labels that identify each curve clearly;
    \item discussion in the main text that states the key physical message of
          the figure, not merely a repetition of the caption.
  \end{itemize}

\item \textbf{Consistency of notation and physical interpretation.}
  The report should define important quantities explicitly and use them
  consistently.  Examples include the density of states $\rho$, the bandwidth
  $\Delta E$, the FGR rate $\Gamma$, the recurrence time $T_\mathrm{rec}$,
  and any derived quantity such as the reorganization-energy measure
  $\Lambda = \sum_k |V_{1k}|^2$ used to organize the $N$-dependence.  If a
  term is used in a nonstandard sense, the report should say so clearly.

\item \textbf{Scientific writing and citation practice.}
  A successful report should read like a coherent scientific argument rather
  than a lab notebook.  This means:
  \begin{itemize}
    \item claims should be supported by equations, figures, or citations;
    \item analytical statements should be distinguished from numerical
          observations;
    \item limitations of the present models should be acknowledged;
    \item literature citations should be included for standard concepts such
          as Fermi's golden rule, Green's-function / self-energy treatments,
          damped two-level (Rabi) dynamics, and possible future generalizations
          such as oscillator-bath models.
  \end{itemize}

\end{enumerate}

\end{document}
